% ===== §6 Semiotics: what the sign can and cannot do; the poem as paradigm =====
\section{What the Sign Can and Cannot Do: The Poem as Paradigm of the Good Circle}
\label{p09:sec:semiotics}

The preceding sections erected a criterion. The good circle is a spiral and not a circle: it returns to its root while accruing an irreducible remainder, a surplus we have named, in the language of geometric phase, a positive holonomy. The sections on epistemology and praxis then asked what finite agents can do in the face of such a criterion, and answered with a disposition rather than a technique---not to \emph{produce} sublimation but to refrain from destroying the curvature under which sublimation arises of itself. Yet that answer left a promissory note unredeemed. For whether we speak of guarding or of repairing, we must \emph{do} something; and everything we can do must pass through a single interface---the \emb{sign}.

This section takes up that interface. It must answer three questions, and the order is dictated by a principle that has governed the series throughout: epistemology before praxis. \emph{Why} can a sign alter the trajectory of a cycle at all (\xref{p09:sec:two-mechanisms})? \emph{When and how} must a sign intervene if it is to bend the cycle toward the good (a question that begins here and is completed only in \xref{p09:sec:core})? And \emph{when and how} does a sign fail (\xref{p09:sec:failures})? The section's summit (\xref{p09:sec:poem}) will then advance what may be among the strongest claims of the paper: that among all signs, the \emb{poem} occupies a singular place---that it is the pure form of the good circle's sign, the Archimedean point of the entire semiotics.

\subsection{The paradox of the sign: the sole operable thing, and the gateway to alienation}
\label{p09:sec:symbol-paradox}

We must begin from a tension, and resist the temptation to dissolve it prematurely; the tension is itself the engine of the section.

On the one hand, \emb{the sign is the only thing we can formally operate, know, and transmit.} The real of love---that connection between two beings which is unsayable, immeasurable, exhausted by no concept---is not, in itself, operable. We cannot ``adjust'' it directly, cannot hold it in the hand and add to or subtract from it. What we can lay hands on is only the sign: a declaration, a gift, a rite, a line set down on paper. All practice must pass through this narrow gate. In this sense the sign is not the ornament of a relation but the sole stratum at which a relation can be deliberately entered. To deny the sign its office is to deny the very possibility of practice.

On the other hand---and this is no less true---\emb{the real of love does not reside in signs.} If a relation subsists on signs alone, it has already stepped into alienation: the signifier slides, multiplies, and is exchanged, yet anchors no real. This has been a constant concern of the series, and it acquires today an unprecedented edge. Technology, and above all the advent of large language models, is driving \emb{the marginal cost of the signifier toward zero}. A machine can generate an inexhaustible supply of endearments, of bespoke declarations, of impeccably calibrated phrasings for a gift. When signs can be produced this cheaply and at this scale, the sign's function as a \emph{credible signal of investment and commitment} collapses. This is not a distant worry but an inflation of the sign already underway.

Set the two faces side by side and the genuine problem of a semiotics comes into view; and it hangs, precisely, upon the geometric framework erected above.

\begin{claim}[The paradox of the sign]
The sign is the sole operable interface upon the dynamics of a relation, yet the sign itself \emph{carries no} geometric phase---phase is accrued along a path, and a path is a process walked through in the real, incompressible. The sign can \emph{initiate, redirect, mark} a cycle; it cannot \emph{substitute} for the cycle's own unfolding. \emph{Alienation} is precisely the error of taking the throwing of a sign for the walking of a path---the attempt to cash out, directly, by means of the sign, a sublimation that ought to have arisen as the by-product of a path actually traversed. And artificial generation has industrialized this very illusion: by rendering ``the throwing of a perfect sign'' all but free, it tempts us, as never before, to counterfeit the real path with the sign.
\end{claim}

This is why ``the sign growing ever cheaper'' is a threat at the core of sustainability---not because the sign has ceased to matter but, on the contrary, because the cheap sign lets the vicious cycle (the cycle of zero area, of immediacy, of the skipped process) masquerade as the good one. A love-letter ghost-written by a machine, flawless in its diction, and a clumsy sentence spoken aloud to one's face, freighted with real hesitation---on the ``quality'' of the sign the former far surpasses the latter; but on the real path each anchors, the latter may carry the whole of the phase while the former tends toward zero. The entire unfolding of this section is a reverberation of this paradox.

\subsection{How a sign alters dynamics: two mechanisms}
\label{p09:sec:two-mechanisms}

To answer ``when should a sign intervene'' one must first answer ``by what power can it intervene at all.'' This is the order of epistemology before praxis: only by grasping the mechanism does practice escape blind trial. I claim that the sign alters the trajectory of a cycle by two \emph{independent} mechanisms, belonging respectively to psychoanalysis and to political economy. They alter the dynamics differently, and must be set out separately before their synthesis can be seen.

\subsubsection{The psychoanalytic mechanism: the sign anchors the topology of the loop}

The first mechanism is structural. In the Lacanian register the sliding of the signifier is in principle endless---meaning defers without rest along the chain, with no natural mooring. A signifying act---an ``I love you,'' a naming, a rite---functions as a \emb{quilting point} (\emph{point de capiton}): retroactively it pins down the meaning that until then floated free, producing a stable signified.\footnote{Lacan's figure is the upholsterer's button that fastens the stuffing: it is the discrete anchors that lend structure to an otherwise loosely sliding fabric. The crux is that the anchoring is \emph{retroactive}---the meaning does not wait, in advance of the sign, to be designated; it is constituted retroactively at the moment the sign falls \parencite{lacan2006}.} 

Restated in the language of dynamics: the sign does not inject energy into the system; it \emph{alters the system's topology}. It manufactures, in the phase space of the relation, a new fixed point or attractor, or---and this is the decisive one for us---it \emb{closes an otherwise open, dissipative trajectory into a loop}. A declaration can bend the course of a relation not, as a rule, because it conveys new information (both parties most often already know), but because it anchors the real diffused between them, as yet unformed, into a structure that can be \emph{repeated, and therefore reproduced}.

This matters for the whole paper, because \emb{without a loop there is no holonomy to speak of}. The geometric phase is a path-integral around a closed loop; and for there to be a closed loop, something must first suture the open trajectory. The sign is that suture. Hence the first epistemic office of the sign:

\begin{claim}[The topological office of the sign]
The sign structures the unformalizable real into a fibre bundle capable of bearing a geometric phase. The sign does not produce sublimation, but it delineates the \emph{loop} within which sublimation can occur. This is the first ground of what the sign \emph{can} do, and the reason it cannot be dispensed with: absent the anchoring of the sign, the real of a relation is merely a dissipative flow, with no remainder to accrue.
\end{claim}

\subsubsection{The political-economic mechanism: the sign coordinates the good equilibrium}

The second mechanism is selective, and a different vocabulary is needed for it. The sustainability of a relation---of any cooperative system---is often a problem of \emb{multiple equilibria}. There exists in the system a ``good'' equilibrium in which the two parties mutually invest and value circulates back to each, and a ``bad'' equilibrium in which they mutually defend and each extracts. Which the system falls into is determined not by who is more virtuous but by \emph{the parties' shared belief about one another's intentions}---common knowledge, in the game-theoretic sense.

The sign---and above all the costly, public, irreversible sign---here functions as a \emb{coordinating device}. A weighty gift, a public rite, a covenant written down and witnessed: their office is to manufacture common knowledge---``we both know, and each knows the other knows, that the game we are playing is the good one.'' It is this layer of shared belief that pushes the system out of the basin of the bad equilibrium and into the basin of the good. This explains a fact long obscured by romantic talk: why the gift must be \emph{costly}, the declaration \emph{public}, the covenant \emph{written and witnessed}---cost and publicity are not ornament but the credibility of the signal, and credibility produces shared belief, and shared belief flips the equilibrium.\footnote{This is the logic of costly signalling \parencite{spence1973,zahavi1975} and of coordination by focal points \parencite{schelling1960}. But placed within the framework of generative justice it acquires a new sense: the political-economic office of the sign lies not in ``conveying information'' but in serving as the focal point that makes the good equilibrium common knowledge.}

\subsubsection{The synthesis of the two mechanisms}

The two mechanisms can now be joined, and the synthesis is itself a new claim. The psychoanalytic mechanism says: the sign \emph{constructs the topology of the loop}, making sublimation structurally possible. The political-economic mechanism says: the sign \emph{selects the basin of the loop}, making the good circle the one jointly chosen. One is topology; the other is selection.

\begin{claim}[The double office of the sign upon dynamics]
The sign alters the dynamics of a relation at two levels: it \emph{delineates} the structure of the phase space (psychoanalysis: anchoring the loop), and it \emph{coordinates}, among multiple stable states, the good branch (political economy: selecting the basin). The sign-inflation of the age of artificial generation attacks both at once---it devalues the anchor (signifiers proliferate and can no longer pin a real), and it strips the signal of its cost (no longer credible, it can coordinate no common knowledge). The crisis of the sign is therefore a \emph{double} crisis.
\end{claim}

Note that the two mechanisms do not reduce to one another. The topological anchoring can take place within a single mind in solitude (a vow one speaks to oneself anchors one's own real), whereas the coordination of an equilibrium is essentially multi-agent. They are two incommensurable faces of the sign's power, and they prefigure the polyphony the paper's final section will avow: one and the same phenomenon shows itself, at different epistemic stations, as different and mutually irreducible truths.

\subsection{How a sign fails: three pathologies}
\label{p09:sec:failures}

Having grasped what the sign \emph{can} do, we can characterize precisely what it cannot. The failure of the sign is not single but takes three structurally distinct modes. To distinguish them is essential, for each answers to a different pathology and demands a different response.

\subsubsection{Inflationary failure: the surfeit of the sign}

The first failure is the \emb{surfeit} of the sign. The signifier multiplies without bound while the real investment of the signified does not grow with it. The contemporary performance of relation---the staged tenderness of social media, the algorithmically suggested phrasing of an anniversary, the machine-generable endearment---is its type. Here both offices of the sign fail at once: anchoring fails (too many signifiers, none can pin a real), and signalling loses faith (cost at zero, no common knowledge can be coordinated). The cycle seems busier, more ``full of content,'' than ever, and yet it idles---its geometric phase tends to zero. This is the dominant pathology of the present, and the direct issue of the sign's inflation.

\subsubsection{Settling failure: the usurpation by the sign}

The second failure is more hidden and more profound---and it is precisely the pivot of this section toward the poem. Here the sign is used to \emph{substitute} for, rather than \emph{invite}, a path: to settle with a sentence a wound that needed time to heal, to discharge with a gift a debt that needed presence to be borne, to ``complete'' with a perfect phrase a process that ought to have been walked. The sign here is no longer the quilting button that anchors the loop; it \emb{declares the loop already closed}---though in fact it is not.

\begin{claim}[Settling failure]
Settling failure is the use of a sign to \emph{cash out} a sublimation never walked through. It counterfeits, as a balance the sign can draw on at once, a phase that ought to have arisen as the by-product of a path. Its essence is fraud---against the other, and against oneself. For phase depends on the path and not on the endpoints: two relations may arrive at ``the same situation,'' yet, having walked different paths, accrue different sublimations. Sublimation cannot be acquired instantaneously; it can only be walked. Every sign that proclaims ``we have arrived'' while in truth it never set out is a settling failure.
\end{claim}

It is this failure that forces an unintuitive insight: \emb{depth admits no shortcut}. In the language of geometric phase, the accrued phase equals the area enclosed by the path---the integral of the curvature around the loop. Hence the phase depends on the path and not on the endpoints. Carried over to relation: what we call ``depth'' just \emph{is} the area the path encloses. One cannot leap over the process and possess a deep relation directly---depth \emph{is} that incompressible area. So ``instant intimacy,'' ``express devotion,'' must be, structurally, vicious cycles: they shrink the area of the loop to the minimum (shortest path, zero waiting), and the geometric phase tends to zero, and no sublimation occurs. The ``immediacy,'' ``acceleration,'' ``efficiency'' extolled by capital's logic are precisely the systematic annihilation of geometric phase---the flattening of a spiral that might have climbed into an oscillation in place. \emb{Slowness, therefore, is a formal necessary condition of the good circle, and not an aesthetic preference.}

\subsubsection{Ossifying failure: the petrification of the sign}

The third failure is \emb{petrification}. A sign that once anchored a good circle hardens into an empty rite that no longer connects to any real---the perfunctory ``I love you,'' the commemoration become mere husk, the gift become duty. The loop is still there, but it is now the pure repetition of zero curvature (what the earlier sections called the ``zombie cycle that refuses to wilt''). The sign here is not at fault; the fault is that it is no longer re-walked in the real; it has decayed from anchor into inertia.

The remedy for ossifying failure echoes, exactly, the ``let it wilt and then regenerate'' of the sense of \emph{wu wei}: to let the old sign die, and let a new anchoring grow back out of the real. To try to rescue a petrified sign by more frequent, more forceful repetition only hastens its rigor mortis---which is the structure recurring throughout (and in the earlier praxis section): the action born of fear manufactures the very outcome it fears.

\subsection{The poem as paradigm of the good circle's sign: the Archimedean point of the semiotics}
\label{p09:sec:poem}

With the three failures before us, a naive conclusion seems to suggest itself: since the sign is so dangerous, let us use fewer signs, return to some presemiotic innocence. This conclusion is false, and instructively so. For \emb{the contrary of settling failure is not ``fewer signs'' but a special kind of sign---one that refuses to settle itself}. And among all signs, the one that brings this ``refusal to settle'' to purity is the \emb{poem}.

This is the core claim I wish to make: the poem is not one example within the semiotics but its \emph{Archimedean point}---a fulcrum from which the whole understanding of ``what a good sign is'' may be levered. Let me unfold, in three layers, the poem's title to being the paradigm of the good circle's sign.

\subsubsection{First layer: the poem structurally refuses closure, and is therefore immune to settling}

A contract, a sentence reading ``the matter is resolved,'' a machine-generated standard endearment---their semiotic function is \emph{termination}: to declare the loop closed, the meaning delivered, the file ready for archiving. The poem cannot do this, and \emph{refuses} to. Every time a poem is read, its meaning unfolds anew; it is never ``read out.''\footnote{Here ``poem'' does not designate the literary genre of lineated text but a manner in which a sign operates---which may equally take place in a passage of prose, a rite, a held gaze, and may fail to take place at all in a formal ``poem'' (see the false poem of \xref{p09:sec:false-poem}).} Precisely because the poem cannot be settled, it cannot be used as the instrument of ``cashing out an unwalked sublimation.'' In other words, \emb{the poem is, in its very form, immune to settling failure.} This is its first title to being a paradigm: it is a self-defending sign, a sign that structurally refuses to be put to usurping use.

\subsubsection{Second layer: the poem moves the site of sublimation from the text to the path of interpretation}

This is the most important layer, and the interface that joins the poem to the geometric framework. The meaning of a poem does not lie in the arrangement of its signifiers---this is why we say ``the formal content of the poem matters less''; the meaning of a poem lies in \emb{its unfolding within the dynamics of the interpreter}.

In the language of geometric phase: the text of the poem is only a loop (or a point of departure) in the base space of situations; meaning, sublimation, is \emph{the holonomy the interpreter accrues by traversing that loop once around}.

\begin{claim}[The meaning of a poem is the holonomy of the interpretive path]
The meaning of a poem $=$ the area the interpreter encloses along its loop. The same poem, for different interpreters---and for the same interpreter at different moments---walks different paths, encloses different areas, accrues different phases. This is why a poem is ``inexhaustible'': its loop permits an area arbitrarily large, and unrepeating, to be enclosed. Each rereading is a spiral and not a circle---it carries home a remainder the last did not.
\end{claim}

From this follows a surprisingly operable measure of ``poeticity'': \emb{the poeticity of a sign equals the supremum of the non-trivial holonomy it can bear.} A slogan, an advertising jingle---its loop is too small, too closed; on a second recitation the phase is already zero; it has been ``read out''; its poeticity tends to zero. A great poem withstands a lifetime of rereading precisely because its loop permits the accrual of an almost unbounded, unrepeating phase. Poeticity is not a vague term of praise but a structural magnitude---a magnitude concerning the capacity of a sign to bear sublimation.

\subsubsection{Third layer: the poem exhibits the epistemic structure of witness}

The third layer fills the gap left by the first mechanism (anchoring). The poem does not transmit to the interpreter a ready-made signified; it \emb{invites the interpreter to walk a real once over in person, and to become, in the walking, a witness.} This is the pure form of that act by which a sign ``anchors the loop without substituting for the path''---the poem, with the fewest signifiers, sets in motion the largest unfolding of the real within the interpreter's own dynamics.

And every sign-intervention done rightly---a declaration, a gift, a covenant---does the very same thing: it does not ``inform'' the other of love; it \emph{invites} the other to walk love once over, in person, within their own dynamics, and there to become a witness. The poem, then, is not one example within the semiotics but \emb{the limit form of all legitimate sign-intervention}---the ``invitation of the path'' brought to purity, the ``settling'' reduced to zero.

\begin{claim}[The poem as the apex of the good sign]
Every good sign-intervention, at the limit, approaches the structure of the poem: it \emph{anchors} the loop while refusing to \emph{close} it, it \emph{initiates} sublimation while refusing to \emph{cash it out}, it \emph{invites} the interpreter to walk the real in person and there to witness generativity itself. Settling failure is precisely the moment a sign betrays its own poeticity and usurps the office of a contract. The crisis of the sign in the age of artificial generation is, in essence, \emph{the poeticity of the sign being devoured by its settleability}---the infinitely generable signifier tempting us to treat every sign as a settleable deliverable, and thereby annihilating holonomy. Hence the practice of resisting alienation just is the \emph{restoration of the sign's poeticity}: to let the sign be, once more, an invitation to a real path, and not the cashing-out of a sublimation.
\end{claim}

This claim carries, besides, a retroactive dividend. The eighth paper of this series treated the lyric's address to the big Other---the poem presents generativity itself, rather than any ready-made structure. There that thesis concerned the lyric; here it is re-sited: the structure of the lyric just \emph{is} the semiotic paradigm of the sustainable good circle. The series, then, is not eight papers and then one more, but a ninth that \emph{retroactively re-anchors} the meaning of the eighth---which is itself a \emph{point de capiton}, a spiral return. \emb{The series enacts, in its form, the content it asserts}: no single paper settles the meaning, which goes on generating new remainders along the path of the papers' mutual rereading.

\subsection{The vicious twin of poeticity: the false poem}
\label{p09:sec:false-poem}

But we must at once set a guard upon this too-beautiful picture, lest it collapse into a romanticization of ``poetic vagueness.'' That the poem refuses to settle is its strength; yet \emb{``never settling'' may slide into another pathology}---an infinite deferral, a game of the signifier that never redeems itself: the pure sliding of signifiers, the cynic's irony, the posture that pleads ``staying open'' as an alibi for never investing. This is in truth a higher form of inflationary failure, wearing the mask of the poem.

Poeticity, then, has its vicious twin. I call it the \emb{false poem}. The difference between the true and the false poem returns, in the end, to geometric phase:

\begin{claim}[The criterion of true and false poem]
The \emph{true} poem's non-closure is in order that sublimation \emph{be lived}---the interpreter walks the path in person and accrues a positive holonomy. The \emph{false} poem's non-closure is in order to \emph{evade} the cost of sublimation---its holonomy is zero, it proclaims ``openness'' in place while no one has truly set out. \emph{Openness} is not, in itself, the good; only the openness that bears a positive phase is. A true poem invites a real path to be walked; a false poem pleads ``openness'' as an alibi for never investing.
\end{claim}

This distinction is crucial, for without it the section would be justly charged with aestheticizing vagueness, with dressing ``nothing is certain'' as profundity. The false poem's refusal of closure is an idling refusal---it never demands of the interpreter the cost of walking a real in person, and so it breeds countless postures and zero sublimation. The true poem's refusal of closure is a generous refusal---it hands the whole weight of sublimation back to the path the interpreter lives. On the surface both ``give no ready answer''; in holonomy they are opposites.

And here the question of the \emph{bearer} of the phase can no longer be deferred: that ``generativity'' the poem presents, that ``existence'' the interpreter witnesses---of what is it the phase? We have twice now arrived at the threshold of one and the same answer---once following the logic of the sign's anchoring, once following the logic of the poem's unfolding---and both independent threads point to a single bearer: \emb{the sublimation of jouissance} \parencite{lacan1998}. The poem is the pure artistic form of this sublimation.\footnote{In the Seminar on \emph{The Ethics of Psychoanalysis} \parencite{lacan1992} Lacan expounds sublimation precisely by way of the relation between the poem and \emph{das Ding}: sublimation is not the satisfaction of desire but the ``raising of an ordinary object to the dignity of the Thing''---encircling that impossible, empty real without filling it. This structure will become, in \xref{p09:sec:core}, the heart of the definition of the appropriate poem.} But to fix this bearer securely, and to say what divides the ``good'' poem from the ``vicious'' one, we must first pass a threshold---the next section shows that poeticity, in itself, \emph{does not guarantee} the good.
